\section{Local inertia and the Euler factors at 3}\label{sec:local-three}
This section computes the arithmetic local representation of the literal
rational descents fixed in Section~\ref{sec:descent}. It combines an explicit
quotient with the local two-section vanishing-cycle calculation and a separate
Frobenius computation. The proof does not use a numerical functional-equation fit.

\subsection{Statement and normalization}
Let $\ell\ne3$, let
\[
 T=\PP^1_z\setminus\{0,\infty,\mu_3\},\qquad
 E_z:y^2=x(x-1)(x-z^3),\qquad
 H_\ell=H^1_!(T_{\overline\Q},\Sym^4 R^1\pi_*\Q_\ell).
\]
Here $H^1_!=\im(H_c^1\to H^1)$, $A$ is pullback by $z\mapsto\zeta_3z$,
and $e_\nt=(2-A-A^2)/3$. The literal inversion $I$ uses
\[
 z\mapsto z^{-1},\qquad (x,y)\mapsto(x/z^3,y/u^3),\quad u^2=z^3.
\]
The square-root sign disappears on the fourth symmetric power. Put
$W_{4,\ell}^{\pm}=(e_\nt H_\ell)^{I=\pm1}$, as in \cite{WalkeV2}.
We use geometric Frobenius, so $\Q_\ell(-1)$ has eigenvalue $3$ at $3$.
The standing descent identities are
\[
 \dim W_{4,\ell}^{\pm}=4,\qquad
 W_{4,\ell}^{-}\simeq W_{4,\ell}^{+}\otimes\chi_{-3}.
\]
They follow from the same finite-cover and involution construction for every
$\ell$, without requiring compatibility at a bad prime.

\begin{theorem}\label{p3:thm:local}
For every $\ell\ne3$, the local quantities of these specified representations are
\[
\begin{array}{c|cccc}
 &\dim(W_{4,\ell}^{\pm})^{I_3}&P_3^{\pm}(T)&\Sw_3(W_{4,\ell}^{\pm})&a_3(W_{4,\ell}^{\pm})\\\hline
 +&1&1-3T&0&3\\
 -&0&1&0&4
\end{array}
\]
where $P_3^{\pm}(T)=\det(1-\operatorname{Frob}_3T\mid
(W_{4,\ell}^{\pm})^{I_3})$.
On $W_{4,\ell}^{+}$, inertia is unipotent with a single Jordan block of length
four; on $W_{4,\ell}^{-}$ it is its ramified quadratic twist.
If $N$ is the nonzero tame monodromy logarithm and $F$ is geometric
Frobenius, then $FNF^{-1}=3^{-1}N$. On the successive one-dimensional
quotients of $0\subset\ker N\subset\ker N^2\subset\ker N^3\subset V$
for $V=W_{4,\ell}^+$, the Frobenius eigenvalues are $3,9,27,81$.
\end{theorem}

In particular, the original $2$-adic representation can be used directly at
$3$. The argument also works at $\ell=5$ and yields the same displayed data.
The dyadic calculation is treated separately below.

\subsection{An exact quotient, including the inversion sign}
Set $t=z+z^{-1}$ and consider
\begin{equation}\label{p3:eq:family}
 C_t:Y^2=X^3+(t-1)(t+2)X^2+(t+2)X.
\end{equation}
Write $a=(t-1)(t+2)$, $b=t+2$. Its invariants are
\begin{align*}
 c_4&=16(a^2-3b)=16(t+2)(t^3-3t-1),\\
 \Delta&=16b^2(a^2-4b)=16(t+2)^3(t-2)(t+1)^2,\\
 j&=\frac{256(t^3-3t-1)^3}{(t-2)(t+1)^2}.
\end{align*}
Over $\Q(z)$ put $d=-z^2/(z+1)$. Direct identities give
\[
 da=-(1+z^3),\qquad d^2b=z^3.
\]
Thus $E_z$ is precisely the quadratic twist of $C_t$ with coefficients
$da,d^2b$. Over $v^2=d$, the map
\[
 \psi_z:E_z\longrightarrow C_t,\qquad(x,y)\longmapsto(x/d,y/v^3)
\]
induces a canonical fourth-symmetric-power isomorphism. Its two choices
differ by $[-1]$, which acts trivially on that symmetric power.

Crucially, $d(1/z)=d(z)/z^3$. Choosing $v'=v/u$ for the inverted parameter gives
\[
 \psi_{1/z}\circ\phi_u(x,y)=(x/d,y/v^3)=\psi_z(x,y).
\]
This proves compatibility with the \emph{literal} $I$ in the manuscript.
The descent is the plus descent; it has no additional constant quadratic twist.

Let $V=\PP^1_t\setminus\{2,-1,-2,\infty\}$ and
$Z=\PP^1_z\setminus\{0,\infty,\mu_3,-1\}$.
Then $Z\to V$ is a finite etale double cover, and the displayed isomorphism
identifies its invariant coefficient system with
$\mathcal F=\Sym^4R^1(C/V)_*\Q_\ell$.
On the compact curves, finite pushforward followed by invariants identifies
the ordinary middle-extension sheaves. This can be checked stalkwise:
the stalk at a missing point is the invariant space of its punctured-trait
monodromy. Finite pushforward and rational finite-group invariants are exact.
Consequently
\begin{equation}\label{p3:eq:plus}
 H^1(\PP^1_{\bar\Q},j_*\mathcal F)\simeq H_\ell^{I=+1}.
\end{equation}
Adding back $z=-1$ on the original side makes no change to its parabolic
cohomology, since its original coefficient sheaf is smooth there.

\subsection{The rank and the extra branch point}
The singular fibres of \eqref{p3:eq:family} in characteristic zero are
\[
\begin{array}{c|rrrr}
 t&2&-1&-2&\infty\\\hline
 \text{type}&I_1&I_2&\mathrm{III}&I_6\\
 \dim\mathcal F^{I_t}&1&1&3&1
\end{array}
\]
At infinity use $r=1/t$, $X'=r^2X$, $Y'=r^3Y$. The resulting coefficients
are $1+r-2r^2$ and $r^3(1+2r)$; $c_4$ is a unit and $\Delta$ has order six.
At $-2$, the orders of $c_4,\Delta$ are $1,3$, giving type III.
Equivalently, the ramified preimage $z=-1$ has elliptic automorphism
$(x,y)\mapsto(-x,iy)$. Its $H^1$ eigenvalues are $i,-i$, so the eigenvalues
of its fourth symmetric power are $1,-1,1,-1,1$.
This explains both the three-dimensional stalk at $-2$ and why that point
must be included in the rank calculation.

The Legendre geometric monodromy is $\mathrm{SL}_2$; finite-index pullback
preserves its connected Zariski closure. The quotient identification therefore
shows that neither $\mathcal F$ nor its dual has a geometric global invariant.
The parabolic Euler-characteristic formula gives
\[
 \dim H^1(\PP^1,j_*\mathcal F)=(4-2)5-(1+1+3+1)=4.
\]
The already defined rank-four $W_{4,\ell}^{+}$ is a subspace of the right
side of \eqref{p3:eq:plus}; equality of dimensions proves
\begin{equation}\label{p3:eq:identification}
 W_{4,\ell}^{+}=H^1_!(V_{\bar\Q},\mathcal F)
\end{equation}
with its full arithmetic action. This step does not use Hodge numbers or a
numerical Frobenius comparison.

\subsection{The collision at 3 and its local cycles}
On $\mathcal X=\PP^1_{\Z_3}$ remove the four sections
$D=\{2,-1,-2,\infty\}$ and let $j:\mathcal U\hookrightarrow\mathcal X$.
The elliptic scheme \eqref{p3:eq:family} is smooth and proper over $\mathcal U$,
since its discriminant is a unit there. Thus $\mathcal F$ is lisse on this
entire open scheme for every $\ell\ne3$.

The sections $t=2,-1$ meet at the special point $x$ with $t=-1\in\F_3$.
Their intersection multiplicity is $v_3(2-(-1))=1$. The other two sections
specialize to $1,\infty$, separately. Near $x$ the invariant $c_4$ is a unit
and the discriminant is a unit times $(t-2)(t+1)^2$.
At $x$ the nodal cubic is
\[
 Y^2=X^3+X^2+X=X(X-1)^2.
\]
The node $(1,0)$ has rational tangents $Y=\pm(X-1)$: the degeneration is
split multiplicative. The two horizontal monodromies on $R^1$ have
multiplicities one and two in the same Tate unipotent direction.

We apply \emph{the local} theorem \cite[\S\S1.6--1.9]{SchollVanishing}, not its
global semistable-family Theorem 2.15. The present family has an additive
point at $-2$, so that global theorem's hypotheses should not be asserted.
For the local theorem, the base curve is smooth, the two sections have
intersection multiplicity one, the coefficient system is lisse off them,
and Tate uniformization makes its local monodromy factor through the two
Kummer characters with multiplicities $1,2$. These are exactly the required
local hypotheses.

Put $B=(R^1\Phi j_!\mathcal F)_x$. Scholl's calculation identifies the
inertia action on $B$ with the rank-five fourth symmetric power of a
nontrivial unipotent rank-two action. Its logarithm is proportional to
$(1+2)N_0=3N_0$, and $B$ is a single unipotent block of length five.
In particular wild inertia acts trivially.
The local boundary sequence \cite[\S1.7]{SchollVanishing} gives
\begin{equation}\label{p3:eq:ABC}
 0\longrightarrow A_x\longrightarrow B\longrightarrow C_x\longrightarrow0,
 \qquad \dim A_x=1,\quad \dim C_x=4.
\end{equation}
Here $A_x$ is the cokernel of the diagonal map from the common invariant
line to the two cusp invariant lines. It is the unique invariant line of $B$.
Thus $C_x$ is a single unipotent block of length four.

\subsection{The global-to-local surjection}
This paragraph supplies the necessary bridge from local cycles to the
specified global representation. It also avoids any assumption about the
absence of other vanishing cycles at the additive point. In particular,
we do not invoke the original deduction of Corollary 2.16 in Scholl's
paper; its published correction is included in the reference
\cite{SchollVanishing}.

On the special open curve $\mathcal U_{\bar\F_3}$,
$H_c^2(\mathcal U_{\bar\F_3},\mathcal F)=0$.
Indeed pull back by the separable double cover $t=z+z^{-1}$.
The exact quadratic-twist identities still hold in characteristic three.
After taking $\Sym^4$, the resulting system is the Legendre system at $z^3$.
The relative Frobenius isogeny from the Legendre curve at $z$ to that at $z^3$
induces an isomorphism on $\ell$-adic $H^1$ for $\ell\ne3$.
The Legendre geometric monodromy is $\mathrm{SL}_2$ in every odd
characteristic \cite[Chapter 15, opening discussion]{Katz}; hence its
fourth symmetric power and its dual have no global invariants. Pullback
detects invariant sections, and Poincare duality gives the stated vanishing.

Apply the nearby/vanishing-cycle triangle to $\mathcal G=j_!\mathcal F$
on the proper curve $\mathcal X$. Since $\mathcal G$ is lisse on the smooth
open, $R\Phi\mathcal G$ is supported on the finite special boundary.
Its degree-one hypercohomology is therefore the direct sum of its
degree-one stalks. The exact sequence contains
\[
 H_c^1(V_{\bar\Q_3},\mathcal F)\longrightarrow
 \bigoplus_{y\in D_{\bar\F_3}}(R^1\Phi\mathcal G)_y
 \longrightarrow H_c^2(\mathcal U_{\bar\F_3},\mathcal F)=0.
\]
Projection gives a surjection onto $B$. Naturality of the boundary maps
and \eqref{p3:eq:ABC} show that the boundary subspace of $H_c^1$ maps into
$A_x$. More explicitly, the two distant cusp stalks have zero contribution
at $x$, while the two colliding invariant lines map through
$(\Q_\ell\oplus\Q_\ell)/\operatorname{diag}\Q_\ell=A_x$ and onto it.
Thus there is an induced
surjection
\[
 H^1_!(V_{\bar\Q_3},\mathcal F)\longrightarrow C_x.
\]
Its source and target have dimension four, so it is an isomorphism.
All constructions are natural over $\Q_3$, and $x$ is rational over $\F_3$;
the isomorphism is equivariant for the full local Galois group.
Together with \eqref{p3:eq:identification}, this proves that $W_{4,\ell}^{+}$
has tame, unipotent inertia with a single length-four block.

\subsection{Arithmetic Frobenius: fixing the sign and the Tate twist}
Scholl's local theorem by itself specifies inertia, not the arithmetic
Frobenius eigenvalue. We give the extra calculation here.
Near the split Tate degeneration, $R^1$ has its natural filtration
\[
 0\longrightarrow\Q_\ell\longrightarrow R^1
 \longrightarrow\Q_\ell(-1)\longrightarrow0.
\]
The fourth symmetric power has five graded pieces
$\Q_\ell(-j)$, $0\le j\le4$.
They are unramified arithmetic coefficient lines extending over the local
base. The split-node tangent calculation removes any torus splitting
character. Alternatively a quadratic splitting character disappears on
the even fourth power.

For a constant unramified coefficient line $L$, the local vanishing cycles
of extension by zero at two colliding, individually rational sections satisfy
\[
 R^1\Phi(j_!L)_x=\operatorname{coker}(L\xrightarrow{\rm diag}L\oplus L)
 \simeq L,\qquad R^q\Phi(j_!L)_x=0\quad(q\ne1).
\]
This follows naturally from the constant-sheaf/closed-boundary exact
sequence and smooth base change. Because both sections are separately
rational, Frobenius does not exchange them. In particular no sign
character and no extra Tate twist enter this formula.

Apply the vanishing-cycle functor to the Tate filtration. The displayed
concentration in degree one makes successive sequences short exact at this
stalk. It yields an arithmetic filtration of $B$ with graded pieces
\[
 \Q_\ell,\ \Q_\ell(-1),\ \Q_\ell(-2),\ \Q_\ell(-3),\ \Q_\ell(-4).
\]
The boundary line $A_x$ is the first piece: it is the difference of the
two split cusp invariant lines. Hence $C_x$ has graded pieces with
$j=1,2,3,4$. Its first one-dimensional filtered piece is inertia invariant;
the single length-four block has just one invariant line. Therefore
\[
 (W_{4,\ell}^{+})^{I_3}\simeq\Q_\ell(-1),\qquad P_3^+(T)=1-3T.
\]
This determines the positive sign arithmetically. It does not infer that
sign from purity, polarization or a numerical functional equation.

The same filtration determines the full ordered monodromy flag. Every
filtered piece is inertia-stable, and a single nilpotent Jordan block
has a unique invariant subspace of each dimension. Thus these pieces
are exactly $\ker N^r$, with graded Frobenius eigenvalues $3^r$ for
$r=1,2,3,4$. Geometric Frobenius acts on the tame inertia parameter by
$3^{-1}$, so taking logarithms gives $FNF^{-1}=3^{-1}N$.

\subsection{The minus descent and the conductor exponents}
The extension $\Q_3(\sqrt{-3})/\Q_3$ is ramified of degree two, hence tame.
Its character $\chi_{-3}$ is nontrivial on inertia and trivial on wild inertia.
On the plus space all inertia operators are unipotent. Choose an inertia
element on which $\chi_{-3}=-1$; its action on the minus space has all
eigenvalues $-1$, so it fixes no nonzero vector. Consequently
\[
 (W_{4,\ell}^{-})^{I_3}=0,\qquad P_3^-(T)=1.
\]
Both spaces are tame, so both Swan conductors vanish. The Artin conductor
formula \cite{Ulmer} is
\[
 a_3(V)=\dim V-\dim V^{I_3}+\Sw_3(V).
\]
It gives $a_3(W^+)=3$, $a_3(W^-)=4$, completing Theorem~\ref{p3:thm:local}.

